\documentclass[12pt,a4paper]{article}

% Paquetes esenciales
\usepackage[utf8]{inputenc}
\usepackage[spanish,es-tabla]{babel}
\usepackage{geometry}
\usepackage{setspace}
\usepackage{amsmath,amssymb}
\usepackage{graphicx}
\usepackage{booktabs}
\usepackage{array}
\usepackage{hyperref}
\usepackage{fancyhdr}
\usepackage{titlesec}
\usepackage{microtype}

% Configuración de Márgenes e Interlineado (Interlineado a 2.0 requerido)
\geometry{top=2.5cm, bottom=2.5cm, left=3.0cm, right=2.5cm}
\doublespacing

% Configuración de Enlaces
\hypersetup{
    colorlinks=true,
    linkcolor=blue!70!black,
    citecolor=blue!70!black,
    urlcolor=blue!70!black
}

% Encabezados y Pies de Página
\pagestyle{fancy}
\fancyhf{}
\rhead{\small\textit{Gestión de Base de Datos - CNV Perú}}
\lhead{\small\textit{Escuela Superior la Pontificia}}
\rfoot{\thepage}

\begin{document}

% ==============================================================================
% PORTADA INSTITUCIONAL
% ==============================================================================
\begin{titlepage}
    \centering
    \vspace*{1cm}
    {\scshape\LARGE Escuela Superior la Pontificia \par}
    \vspace{0.8cm}
    {\scshape\Large Carrera Profesional de Ingeniería de Sistemas de Información \par}
    \vspace{0.3cm}
    {\large Ciclo: VIII -- B \par}
    \vspace{1.5cm}
    
    {\huge\bfseries GESTIÓN DE BASE DE DATOS Y BIG DATA CON DATOS ABIERTOS REALES\par}
    \vspace{0.5cm}
    {\Large\bfseries Arquitectura Dimensional, Normalización 3FN, Pipeline ETL, Modelo Predictivo y Dashboard del Certificado de Nacido Vivo en el Perú (2015 -- 2025)\par}
    \vspace{2.0cm}
    
    \begin{flushleft}
        \large
        \textbf{Curso:} Gestión de Base de Datos / Big Data \\
        \textbf{Docente:} Ing. Erick Jhonatan Palomino Ayala \\
        \textbf{Integrantes:} \\
        \hspace{1cm} $\bullet$ Curahua Morales, Any Miluska
    \end{flushleft}
    
    \vfill
    {\large Ayacucho, Perú \par}
    {\large 2026 \par}
\end{titlepage}

\newpage
\tableofcontents
\newpage

% ==============================================================================
% SECCIÓN 1: INTRODUCCIÓN
% ==============================================================================
\section{Introducción}

La gestión estratégica de grandes volúmenes de datos (\textit{Big Data}) constituye una herramienta fundamental para el diagnóstico y la toma de decisiones en el sector salud. En el Perú, la Plataforma Nacional de Datos Abiertos pone a disposición pública los registros del Certificado de Nacido Vivo (CNV), administrado por el Ministerio de Salud (MINSA) y el Registro Nacional de Identificación y Estado Civil (RENIEC).

El presente informe documenta el desarrollo técnico de una solución integral de ingeniería de datos aplicada sobre \textbf{4,904,793 registros} correspondientes al período \textbf{2015 -- 2025}. Se aborda el modelado relacional en Esquema Estrella, la normalización hasta la Tercera Forma Normal (3FN), el desarrollo de scripts T-SQL sin sentencias \texttt{GO}, el pipeline ETL en Python, el entrenamiento de un modelo predictivo por regresión lineal y el despliegue de un Dashboard Web interactivo en Vanilla JavaScript.

% ==============================================================================
% SECCIÓN 2: OBJETIVOS
% ==============================================================================
\section{Objetivos}

\subsection{Objetivo General}
Diseñar e implementar una arquitectura de base de datos relacional normalizada y analítica para procesar, modelar, predecir y visualizar los nacimientos en el Perú entre 2015 y 2025.

\subsection{Objetivos Específicos}
\begin{enumerate}
    \item Demostrar cuantitativamente la inviabilidad del análisis en hojas de cálculo tradicionales (Microsoft Excel) frente a un volumen de 4.9 millones de filas.
    \item Construir un Esquema Estrella normalizado en 3FN que reduzca la redundancia y optimice el almacenamiento.
    \item Desarrollar los scripts T-SQL en Microsoft SQL Server con restricciones de integridad, índices cubrientes, vistas y 20 procedimientos almacenados estandarizados.
    \item Implementar un pipeline ETL en Python con la librería Pandas para la limpieza, imputación y estandarización de códigos de Ubigeo INEI.
    \item Entrenar un modelo de regresión lineal para proyectar la natalidad y la tasa de cesáreas hacia el quinquenio 2026 -- 2030.
    \item Construir un Dashboard Web interactivo en HTML5, CSS y JavaScript Vanilla que consuma los datos limpios en formato JSON.
\end{enumerate}

% ==============================================================================
% SECCIÓN 3: FUENTE OFICIAL DE LOS DATOS
% ==============================================================================
\section{Fuente Oficial de los Datos}

\begin{itemize}
    \item \textbf{Entidad Productora:} Ministerio de Salud del Perú (MINSA) / RENIEC / Oficina de Tecnologías de la Información (OTI).
    \item \textbf{Plataforma:} Plataforma Nacional de Datos Abiertos del Perú (\url{https://www.datosabiertos.gob.pe/}).
    \item \textbf{Nombre del Dataset:} \textit{Registros de Nacidos Vivos en el Perú 2026 (Certificado de Nacido Vivo - CNV)}.
    \item \textbf{Período Analizado:} 11 años completos (2015 al 2025).
    \item \textbf{Volumen:} 4,904,793 filas, 22 variables y 791.35 MB de tamaño físico.
\end{itemize}

% ==============================================================================
% SECCIÓN 4: DESCRIPCIÓN DEL DATASET
% ==============================================================================
\section{Descripción del Dataset}

El dataset oficial contiene 4,904,793 registros. Una hoja de cálculo convencional de Microsoft Excel posee un límite máximo de 1,048,576 filas. Por consiguiente, el dataset excede en \textbf{367.8\% (4.68 veces)} la capacidad física de Excel, provocando truncamiento crítico de 3.85 millones de registros y justificando la implementación de motores de bases de datos relacionales optimizados.

Las 22 variables comprenden métricas biométricas neonatales (\texttt{PESO\_NACIDO}, \texttt{TALLA\_NACIDO}, \texttt{DUR\_EMB\_PARTO}, \texttt{sexo\_nacido}), características maternas (\texttt{Edad\_Madre}, \texttt{Estado\_Civil}, \texttt{Nivel\_Intrucción\_Madre}, \texttt{DESC\_OCUPACION}) e información asistencial (\texttt{IdUbigeoInei}, \texttt{Ipress}, \texttt{Condicion\_Parto}, \texttt{Atiende\_Parto}, \texttt{Financiador\_Parto}).

% ==============================================================================
% SECCIÓN 5: MODELO ENTIDAD-RELACIÓN Y ARQUITECTURA DIMENSIONAL
% ==============================================================================
\section{Modelo Entidad-Relación y Arquitectura Dimensional}

Se implementó un Esquema Estrella (\textit{Star Schema}) con una tabla principal de hechos y seis tablas dimensionales:
\begin{itemize}
    \item \textbf{Tabla de Hechos:} \texttt{FACT\_NACIMIENTO} (almacena el evento del parto, claves foráneas, métricas continuas y flags booleanos: \texttt{es\_bajo\_peso}, \texttt{es\_prematuro}, \texttt{es\_madre\_adolescente}, \texttt{es\_cesarea}).
    \item \textbf{Dimensiones Maestras:} \texttt{DIM\_TIEMPO}, \texttt{DIM\_UBIGEO}, \texttt{DIM\_MADRE\_PERFIL}, \texttt{DIM\_CONDICION\_PARTO}, \texttt{DIM\_ATENCION\_SALUD}, y \texttt{DIM\_IPRESS}.
\end{itemize}

% ==============================================================================
% SECCIÓN 6: PROCESO DE NORMALIZACIÓN
% ==============================================================================
\section{Proceso de Normalización (1FN, 2FN y 3FN)}

El proceso formal se ejecutó en tres etapas:
\begin{enumerate}
    \item \textbf{1FN:} Eliminación de atributos multivaluados, asegurando atomicidad y definiendo la clave primaria \texttt{id\_nacimiento BIGINT}.
    \item \textbf{2FN:} Erradicación de dependencias parciales separando catálogos geográficos e institucionales.
    \item \textbf{3FN:} Eliminación de dependencias transitivas en jerarquías distritales y sociodemográficas.
\end{enumerate}
La normalización redujo el tamaño de la base de datos de 791.35 MB a aproximadamente 281.70 MB, logrando un \textbf{ahorro del 64.4\% en disco}.

% ==============================================================================
% SECCIÓN 7: SCRIPTS SQL EN MICROSOFT SQL SERVER
% ==============================================================================
\section{Scripts SQL en Microsoft SQL Server}

Todos los scripts fueron desarrollados en T-SQL sin sentencias \texttt{GO}, utilizando terminadores con punto y coma:
\begin{itemize}
    \item \texttt{00\_ejecutar\_todo.sql}: Script maestro de instalación en un solo bloque.
    \item \texttt{04\_01\_creacion\_bd.sql} a \texttt{04\_07\_vistas.sql}: Creación de base de datos, tablas, PKs, FKs, muestra de datos, índices cubrientes y vistas.
    \item \texttt{04\_08\_procedimiento\_almacenado.sql}: Colección de 20 procedimientos almacenados estandarizados (\texttt{sp\_MostrarTodosNacimientos}, \texttt{sp\_NacimientosPorRegion}, \texttt{sp\_VerificarAlertasNeonatales}, \texttt{sp\_TablaTemporalEstadisticas}, \texttt{sp\_RecorrerDepartamentosCursor}, etc.).
\end{itemize}

% ==============================================================================
% SECCIÓN 8: PROCESO ETL EN PYTHON
% ==============================================================================
\section{Proceso ETL en Python}

El script \texttt{05\_etl\_proceso.py} realiza la extracción de archivos Excel (\texttt{pd.read\_excel}) y CSV, imputa valores nulos con la mediana nacional (peso: 3,250 g, gestación: 39 semanas), formatea códigos de Ubigeo a 6 dígitos con ceros a la izquierda (\texttt{zfill(6)}) y carga lotes optimizados a SQL Server.

% ==============================================================================
% SECCIÓN 9: ANÁLISIS DE PATRONES Y SECUENCIAS TEMPORALES
% ==============================================================================
\section{Análisis de Patrones y Secuencias Temporales}

\subsection{Evolución Multianual}
Entre 2015 y 2018 se registró un incremento sostenido de 417,368 a 493,990 partos (+18.36\%). A partir de 2019 se inicia una contracción continua, agudizada por la pandemia en 2020 (-4.85\%), alcanzando 376,786 nacimientos en 2025 (-23.73\% acumulado).

\subsection{Estacionalidad Mensual}
Se identifica un patrón estacional bimodal recurrente con picos en marzo (436,279) y mayo (421,288), y un mínimo en noviembre (384,140).

\subsection{Concentración Regional}
Lima concentra el 29.58\% (1,451,019 nacimientos), seguida de Piura (5.84\%) y La Libertad (5.52\%).

% ==============================================================================
% SECCIÓN 10: MODELO PREDICTIVO DE MACHINE LEARNING
% ==============================================================================
\section{Modelo Predictivo de Machine Learning}

Se formuló un modelo de Regresión Lineal Simple (\texttt{07\_modelo\_predictivo.py}):

\begin{equation}
    \hat{Y}_{\text{nacimientos}} = -6,998.95 \cdot X + 14,583,724.73 \quad (R^2 = 0.3338, \text{ RMSE} = 31,266.25)
\end{equation}

\begin{equation}
    \hat{Y}_{\text{cesareas (\%)}} = +0.4091 \cdot X - 788.03 \quad (R^2 = 0.8814)
\end{equation}

\begin{table}[htbp]
    \centering
    \caption{Proyección Quinquenal de Indicadores (2026 -- 2030)}
    \begin{tabular}{cccc}
        \toprule
        \textbf{Año ($X$)} & \textbf{Nacimientos Estimados} & \textbf{Tasa Cesáreas (\%)} & \textbf{Tendencia} \\
        \midrule
        2026 & 403,861 & 40.79\% & Decreciente \\
        2027 & 396,862 & 41.20\% & Decreciente \\
        2028 & 389,863 & 41.61\% & Decreciente \\
        2029 & 382,864 & 42.02\% & Decreciente \\
        2030 & 375,865 & 42.43\% & Decreciente \\
        \bottomrule
    \end{tabular}
\end{table}

% ==============================================================================
% SECCIÓN 11: GRÁFICOS Y VISUALIZACIONES DEL DASHBOARD WEB
% ==============================================================================
\section{Dashboard Web Interactivo}

Desplegado en el directorio \texttt{/dashboard} (\texttt{index.html}, \texttt{styles.css}, \texttt{script.js}, \texttt{datos.json}), implementado en HTML5, CSS3 y JavaScript puro sin frameworks. Proporciona:
\begin{itemize}
    \item Carga asíncrona reactiva vía \texttt{fetch('datos.json')}.
    \item Filtros interactivos en tiempo real por Departamento, Región Natural y Año.
    \item Gráficos dinámicos con Chart.js (serie histórica sólida + proyección punteada, dona de vía de parto, barras departamentales y estacionalidad).
    \item Matriz regional interactiva con ordenamiento por columnas y búsqueda por texto.
\end{itemize}

% ==============================================================================
% SECCIÓN 12: INTERPRETACIÓN DE RESULTADOS
% ==============================================================================
\section{Interpretación de Resultados}

\begin{enumerate}
    \item \textbf{Transición Demográfica:} La caída sostenida de la natalidad exige planificar la oferta educativa básica y prever el envejecimiento poblacional en el mediano plazo.
    \item \textbf{Sobreutilización de Cesáreas:} La tasa nacional de 38.47\% (superando el 49\% en Tumbes, Callao y Lima) sobrepasa el límite internacional de la OMS (15\%), evidenciando sobrecostos médicos y la urgencia de auditorías clínicas.
    \item \textbf{Vulnerabilidad Social:} Regiones como Huancavelica (10.20\%) y Loreto (9.80\%) presentan la mayor prevalencia de bajo peso al nacer, vinculada a la falta de controles prenatales oportunos.
\end{enumerate}

% ==============================================================================
% SECCIÓN 13: CONCLUSIONES
% ==============================================================================
\section{Conclusiones}

\begin{enumerate}
    \item Se demostró la necesidad técnica de herramientas de Big Data para procesar 4.9 millones de filas que desbordan las hojas de cálculo tradicionales.
    \item El Esquema Estrella en 3FN redujo el tamaño de almacenamiento en un 64.4\% y aceleró las consultas OLAP a tiempos sub-segundo.
    \item La colección de 20 procedimientos almacenados en T-SQL sin sentencias \texttt{GO} garantiza robustez y estandarización académica.
    \item El modelo predictivo y el Dashboard Web en Vanilla JS proporcionan una herramienta de inteligencia sanitaria para la toma de decisiones informadas en el MINSA y los Gobiernos Regionales.
\end{enumerate}

\end{document}
